% Generated by benchmark/export_tex.py from benchmark/runs/locomo-final; do not edit by hand.
\begin{table}[!htbp]
\caption{Published LoCoMo results \cite{chhikara2025mem0} beside this work (GPT-4o-mini judge)}
\label{tab:bench-published}
\centering
\begin{reporttable}
\begin{tabularx}{\textwidth}{|Y|r|r|r|}
\hline
\textbf{System} & \textbf{J (\%)} & \textbf{Context tokens} & \textbf{J minus own full context} \\
\hline
Full-context & 72.90 & 26,031 & +0.00 \\ \hline
Mem0g & 68.44 & 3,616 & -4.46 \\ \hline
Mem0 & 66.88 & 1,764 & -6.02 \\ \hline
Zep & 65.99 & 3,911 & -6.91 \\ \hline
RAG (best, k=2, 256-token chunks) & 60.97 & -- & -11.93 \\ \hline
LangMem & 58.10 & 127 & -14.80 \\ \hline
OpenAI & 52.90 & 4,437 & -20.00 \\ \hline
A-Mem & 48.38 & 2,520 & -24.52 \\ \hline
\textbf{This work: whole conversation (Haiku 5.5)} & 88.83 & 22,400 & +0.00 \\ \hline
\textbf{This work: memory engine (Haiku 5.5)} & 83.18 & 6,790 & -5.65 \\ \hline
\textbf{This work: plain retrieval (Haiku 5.5)} & 80.06 & 4,473 & -8.77 \\ \hline
\end{tabularx}
\end{reporttable}
\par\smallskip{\footnotesize Published systems answer with GPT-4o-mini; this work answers with Claude Haiku 5.5, so absolute scores are not comparable. The last column compares each system with the full-context baseline of its own study. Context tokens for published systems are from the Mem0 paper.}
\end{table}
